\section{Routes that do not work, with their numbers}\label{app:routes}

Sections~\ref{sec:mfg} and~\ref{sec:inc} both reached their form by elimination, and the eliminated
routes are worth recording with the numbers that killed them: each is a natural thing to try, and each
fails for a reason that says something about the model. The first group belongs to the increment
argument of Section~\ref{sec:inc}, where the allowance to be met is $0.0661$ at the newborn stage; the
second reports the measurements behind the closure of the mean field route in Section~\ref{sec:mfg}.

\subsection*{The increment route}

\subsubsection*{Where the failure lives}

Measured at $\gamma=5$ over $r\in[4\%,6\%]$, the failure reaches down to $a=0$ --- so no threshold in
assets confines it, and a Markov bound on a thin upper tail is worthless: the sup of the failure
below any threshold is the global sup, and the resulting bound is the uniform one that already
failed. But sorted by earnings state the picture is different. With Aiyagari's seven-state chain the
sup of the failure over assets, state by state, is
\[
  (0,\;0,\;0,\;0,\;0,\;0.0143,\;0.0460),
\]
zero in the bottom five states and confined to the top two, whose invariant mass is
$\nu_6+\nu_7=0.095$. At $\gamma=8$ it reaches the fourth state and the profile is
$(0,0,0,0.0014,0.0118,0.0333,0.0701)$. The failure of Theorem 1 is an \emph{earnings-state}
phenomenon: it is the high-earners, who are few.

\subsubsection*{The minimum-based step amplifies}

The Theorem-1 step can be made quantitative without touching a power mean. Suppose the failure
$D=c_2-c_1-\Delta r\,a$ is positive at $(a,z)$. Then $b_1>b_2\ge0$, so the low-rate Euler equation
holds above the floor and the high-rate one under the cap; the stage-$k$ failure bound and
monotonicity of consumption in assets give $C_2(z',b_2)\le C_1(z',b_1)+\Delta r\,b_1+\varepsilon_k(z')$,
and $(C+s)^{-\gamma}=C^{-\gamma}(1+s/C)^{-\gamma}\ge C^{-\gamma}(1+\sigma)^{-\gamma}$ with
$\sigma=\max_{z'}s(z')/C_1(z',b_1)$. Chaining the two Euler relations,
\begin{equation}\label{eq:qstep}
  D\;\le\;c_1\Bigl[(R_1/R_2)^{1/\gamma}(1+\sigma)-1\Bigr]-\Delta r\,a .
  \tag{Q}
\end{equation}
This is a genuine quantitative step: it recovers the no-failure criterion when the bracket is
nonpositive, and bounds the failure otherwise. It also amplifies. The derivative of \eqref{eq:qstep}
in $\varepsilon_k$ is $c_1(R_1/R_2)^{1/\gamma}/\min_{z'}C_1(z',b_1)$, measured at $4.7$ at
$\gamma=5$ and $4.2$ at $\gamma=8$. Over sixty stages the recursion started from the exact terminal
condition $\varepsilon_0=0$ reaches $V\approx2.8\times10^{3}$ at $\gamma=5$ against a true supremum of
$0.046$ --- and with the sandwich floor in place of the true $C_1$, $10^{60}$. This is the
amplification the rate-response route already ran into, with a per-step multiplier bounded below by
$1/\Thorn>1$; the life cycle does not escape it, it merely makes it finite.

\subsubsection*{The direct bound is too loose}

The alternative is to bound the failure at each stage separately, with no induction, from the
sandwich: $c_2\le\kappa^2_{k}(m_2+H^2_{k})$ and $c_1\ge\kappa^1_{k}m_1$. That gives $1.835$ against
an allowance of $0.032$, too loose by a factor of $58$, and the source of the looseness is visible:
the sandwich slack is of order $\kappa_kH_k$, measured at $0.96$, an order-one quantity where the
target is of order $0.03$.

\subsubsection*{What the two failures have in common}

The remaining gap is not a matter of assembling existing bounds. It needs an estimate of the
rate-response of consumption whose error is small \emph{relative to the response itself} --- an
absolute error below $0.03$ on a quantity of size about one --- and no sandwich in this programme is
of that quality. That is a sharper statement of the requirement than the programme has had before,
and it is attached to a single distribution-free object rather than to Theorem 1 in full.

\subsubsection*{Bounding the multiplier}

$A$ has to be bounded from primitives, and the obvious bound destroys the gain. Termwise,
$\sum w\,x^{-\gamma-1}\le\sum w\,x^{-\gamma}/\min x$, so $A\le M/\min x$
(\lean{sum\_mul\_rpow\_pred\_le}) --- and $M/\min x$ is precisely the multiplier that bounding the
expectation by its smallest term gives. Measured, $A$ runs from $1.04$ to $1.19$ while $M/\min x$ runs
from $4.7$ to $5.9$; substituting the bound raises $m_k$ to $2.4$--$2.6$ and the product to $10^{20}$.
The crude route returns exactly to where it started.

So $A$'s smallness is a statement about the \emph{weights} of the earnings chain and not its support:
with $\rho=0.9$ the transition puts almost all mass on neighbouring states, where next period's
consumption is close to its certainty equivalent, while the far states that make $\min x$ small carry
almost no probability.

\subsubsection*{Near states and far states}

The obvious way to use the weights is to split the states, bounding the near ones by proximity and
the far ones by their probability. It fails, and instructively. With $d=1$ the far mass is
$3\times10^{-3}$, so the far term is negligible --- and the bound is still $2.48$ against an
allowance of $1.066$, because the near term takes a \emph{maximum} over neighbouring states and
$(C/M)^{-\gamma-1}$ reaches $2.5$ even there. Widening to $d=2$ makes it worse, at $4.99$. The
averaging is the whole of the effect, and any bound that replaces it by a maximum on a subset throws
it away.

\subsubsection*{Why level brackets cannot bound a variance}

Two ways to bound $\operatorname{Var}_w(u)\le0.066$ were tried. Both fail, and the second fails for a
reason worth stating.

Bounding it through \emph{levels} --- $\mathbb E[u^2]=M^2\sum w/C^2$ with the ceiling on $M$ and a
floor on $C$ --- gives $5.5$ with the constraint-incidence floor and $11.0$ with the lower sandwich,
against an allowance of $0.066$: too loose by factors of $84$ and $166$. The reason is the same one
that closed the two-sandwich route at $\gamma=2$, one level down. A variance is a measure of
\emph{spread}, and a per-state bracket of width comparable to the spread contributes a spurious
spread larger than the real one; squaring makes it worse.

Bounding the spread directly is the right instinct. Across income states at fixed assets, cash on
hand differs by $y(z')-y(z'')$, so the propensity brackets the consumption difference:
$\kappa_k|\Delta y|\le|\Delta C|\le s|\Delta y|$ with $s$ any proved upper bound on the propensity.
With the trivial $s=1$ and the incidence floor anchoring the level, maximising
$\operatorname{Var}_w(M/C)$ over the admissible set --- $C(1)$ in its bracket, each increment at
either end --- gives $0.393$ at the newborn stage, against the allowance of $0.066$. Short by a
factor of six. (A single admissible profile gives $0.044$ and appears to pass; it is not the worst
case, and the worst case is what a bound must survive.)

Scaling suggests the slope bound needed is about $s\le0.41$. That cannot be uniform: the true
cross-state slope is $\kappa_k\approx0.04$ where the household is unconstrained but exactly $1$ where
the borrowing constraint binds.

\subsubsection*{The slope, state by state}

The split is available. The saving floor of \lean{le\_stagePolicy\_state\_on} certifies that the
constraint cannot bind at $(b,z')$ whenever $(1-\kappa_k)(y(z')+Rb)>\kappa_kH_k(z')$, a primitive
test, and it certifies $89\%$ to $94\%$ of the pairs that arise. Only $1.6\%$ are in fact
constrained.

What it does not do is let the constrained states be neglected. Measured, there are pairs $(z,b)$ at
which the states carrying the entire transition mass are constrained --- low assets and low earnings,
where the household consumes its cash on hand --- so the weights cannot be used to make that regime
disappear.

The way through is not a better slope bound on that regime but a different handle on it. Where the
constraint binds, consumption is cash on hand \emph{exactly}: $C(z')=y(z')+Rb$, with no bracket at
all, so the spread of $C$ across those states is exactly the spread of $y$, which is known. The two
regimes want treating differently in kind --- an identity on the constrained set, a slope bound on
the unconstrained one --- rather than a single slope bound made state-dependent. The worst-case
worst-case maximisation over the bracket is loose precisely because it gives the constrained states a
bracket they do not have.

\subsubsection*{Branch and bound with an inconsistent bound}

The natural certificate is branch and bound: on a sub-box $[c^-,c^+]$ the harmonic mean is increasing
in each coordinate, so $M\le M(c^+)$ and $u_i\le M(c^+)/c^-_i$, giving an upper bound on
$\operatorname{Var}_w(u)$; branch, prune the slope-infeasible boxes, and refine. It does not converge:
with four thousand nodes per state the certified bound is $4.34$ against an allowance of $0.066$, and
the worst node is the lowest earnings state, where the bracket is widest.

The bound is the problem, not the budget. It evaluates the mean at the top of the box and the
individual consumptions at the bottom, which no single point achieves --- an inconsistency that no
amount of branching removes, because it is reintroduced at every level.

\subsubsection*{Incomplete searches, and the profile they miss}

Enumerating those endpoints settles the question, and not in favour of the bracket. A forward sweep
over the chain --- each coordinate at an endpoint of the range its predecessor leaves --- gives
$0.0414$, and random probes never beat it, which looks like closure. Enumerating patterns in both
directions, so that a coordinate may also be pinned against its successor, gives $\mathbf{0.0874}$ at
the newborn stage, against an allowance of $0.0661$. The admissible set provably contains a profile
that violates the requirement. The forward sweep was incomplete, and so were the coordinate ascent and
the random search that preceded it, all three by a factor of two.

\subsubsection*{A first certificate, in floating point}

The upper bound follows from the same unimodality that justified enumerating. On \emph{any} box the
objective is maximised at a box vertex, so the maximum over a box is \emph{exactly} the maximum over
its $2^{7}$ vertices --- an exact and consistent bound, unlike the one that defeated branch and bound
before, which put the mean at the top of the box and the consumptions at the bottom. Branch and bound
with it: discard a box when it is disjoint from the polytope, by an interval test on the slope and
concavity constraints, or when its exact maximum is at most the allowance; otherwise split the widest
coordinate. If the stack empties, every point of the polytope lay in a discarded box and the maximum is
certified below the allowance.

It empties. At the newborn stage, over forty asset values in each of the seven earnings states, every
node is discarded, with at most $6971$ boxes examined at the worst of them.

The newborn stage is the one that matters. The allowance is $\Thorn/\bigl(q(1-\kappa_kR)\bigr)-1$ and
$\kappa_k$ rises as the horizon shortens, so the allowance rises too --- $0.0661$ at $k=59$, $0.0725$
at $k=40$, $0.0969$ at $k=20$, $0.2586$ at $k=5$ --- while the variance stays at $0.0415$ throughout.
Certifying the youngest cohort certifies every age, and that step is an argument rather than a sample.

\subsubsection*{Its two gaps}

At the calibration, at logarithmic utility, over $r\in[4\%,6\%]$: every input of the chain holds. The
error carried between stages is damped (Theorem~\ref{thm:mpcfloor}, \lean{damped\_le\_one}), the
gradient sum is one plus a variance (Proposition~\ref{prop:gradvar}), and the variance is certified
below what the damping allows. Theorem 1 then holds at \emph{every} state, hence capital supply rises
with the rate, hence the stationary equilibrium is unique.

\subsection*{The mean field condition, measured}

Section~\ref{sec:mfg} rests on two measurements of the Lasry--Lions pairing and two more of the
condition \eqref{eq:W} it reduces to. They are here.

\subsubsection*{Date by date it fails, and it fails at birth}

Take two economies differing only in $\beta$, each at its own equilibrium, so that each
distribution is priced by the aggregate it actually carries. Since $\beta$ does not appear in
$F_t$, this is a legitimate test of the date-$t$ condition. At the calibration of
Section~\ref{sec:eqm} ($\alpha=9/25$, $\delta=2/25$, $J=60$, Aiyagari's earnings chain), with
$\beta=0.960$ against $\beta=0.950$:

\begin{center}
\begin{tabular}{lrrrrr}
\toprule
$\gamma$ & $0.5$ & $1$ & $2$ & $3$ & $5$\\
\midrule
dates with $P_t<0$ (of $60$) & $3$ & $4$ & $5$ & $6$ & $6$\\
$P_0$ & $-0.00147$ & $-0.00227$ & $-0.00324$ & $-0.00419$ & $-0.00804$\\
$\sum_t P_t$ (undiscounted) & $+0.271$ & $+0.175$ & $+0.091$ & $+0.061$ & $+0.047$\\
\bottomrule
\end{tabular}
\end{center}

The failures are at the start of life, and they survive as the two equilibria close up:
$P_0/(\Delta\beta)^2$ converges to $-25.6$ at $\gamma=1$ and $-44.4$ at $\gamma=3$ as
$\Delta\beta$ falls from $0.02$ to $0.00125$, with the count of negative dates unchanged
throughout. So this is a property of the derivative, not an artefact of a wide gap.

Date zero is the sharp case, and it is worth spelling out, because it is exactly the common
initial distribution that was supposed to rescue the argument. Every cohort is born with nothing,
so $a_0=0$ for all of the mass in both economies, and the affine-utility pairing above is
identically zero: with no assets to price, the aggregate coupling has no grip at all. But the wage
still differs, and the utility gain from the higher wage is weighted by marginal utility, which is
higher in the economy that saves more. At $\gamma=1$ mean date-zero saving is $0.2937$ against
$0.2736$, and $P_0$ is strictly negative. The common start does not make the date-zero pairing
vanish; it removes the one term that could have made it positive.

\subsubsection*{Summed, it holds where the equilibrium lives}

The summed condition is a different matter. Test it over the family the equilibria belong to:
$m(\hat K)$ generated by optimal life-cycle behaviour at the prices of $\hat K$, priced inside $F$
by the aggregate $A(\hat K)$ it actually carries. Every equilibrium is an $m(\hat K)$ with
$A(\hat K)=\hat K$, so monotonicity over pairs from this family is exactly what
Theorem~\ref{thm:mfgagg} needs.

\begin{center}
\begin{tabular}{lrrrr}
\toprule
$\gamma$ & $1$ & $2$ & $3$ & $5$\\
\midrule
pairs with $\sum_t\beta^tP_t<0$, $\hat K/K^*\in[0.5,2]$ & $1/66$ & $0/66$ & $0/66$ & $0/66$\\
pairs failing inside $K\in[3.55,7.40]$ & $0/36$ & $0/36$ & $0/36$ & $0/28$\\
least margin $\sum_t\beta^tP_t/(A_1-A_2)^2$ there & $+0.075$ & $+0.097$ & $+0.121$ & $+0.338$\\
\bottomrule
\end{tabular}
\end{center}

The bracket $K\in[3.55,7.40]$ is the one $r\in[2\%,8\%]$ of Theorem~\ref{thm:exist5}, inside which
the equilibrium is proved to lie. Within it the summed pairing is strictly positive for every pair
and every risk aversion tested, with a margin bounded away from zero, and the margin grows with
$\gamma$. The single failure is at $\gamma=1$ between two candidates one of which sits outside the
bracket entirely.

\subsubsection*{A pointwise condition}

Submodularity of the flow $u\bigl(w(K)y+(1+r(K))a-a'\bigr)$ in $(K,a)$ is, for constant relative
risk aversion, exactly
\begin{equation}\label{eq:P}
  \gamma\,(1+r(K))\,(a_t-Ky_t)\;\le\;c_t \qquad\text{at every age and state,}
  \tag{P}
\end{equation}
which implies \eqref{eq:W} and hence the summed condition. It is free wherever the household holds
less than its share of the capital stock, and it is hopeless where it does not: at the calibration
it is violated on $27\%$ to $45\%$ of the mass, at every risk aversion. Requiring each date to
carry itself is precisely what the summed condition was introduced to avoid, and \eqref{eq:P} is
that mistake made again one level down.

\subsubsection*{Where the pointwise ordering fails}

\eqref{eq:W} is a single scalar comparison, and it is weaker than the pointwise ordering of the two
plans' asset paths. Measured at the calibration, on plans drawn from the candidate family at
$\hat K/K^*=0.90$ against $1.10$ --- so the two plans face different wages as well as different
rates, which is the comparison \eqref{eq:W} itself involves:

\begin{center}
\begin{tabular}{lrrrr}
\toprule
$\gamma$ & $1$ & $2$ & $3$ & $5$\\
\midrule
$\Phi_1-\Phi_2$ at the midpoint $K$ & $+39.36$ & $+12.67$ & $+7.34$ & $+7.32$\\
mass on which $a^1_t\ge a^2_t$ FAILS & $0.0000$ & $0.0025$ & $0.0409$ & $0.2006$\\
the same, at a common wage & $0$ & $0$ & $0$ & $0.0091$\\
\bottomrule
\end{tabular}
\end{center}

The second row is the one to read against Light's Theorem 1, which is a statement about the
household at a \emph{fixed} wage: in wage units the policy ordering is exact --- violated on zero
mass, not on small mass --- for $\gamma\le3$, and fails on under one percent of the mass at
$\gamma=5$ (on $6.7\%$ at $\gamma=8$). The first row, which lets the wage move along the frontier as
well, overstates the failure of Theorem 1 by more than an order of magnitude; it is the right
measure for \eqref{eq:W} and the wrong one for Theorem 1. So \eqref{eq:W} is weaker than the
pointwise ordering and does survive where that ordering fails, but the margin between them at
$\gamma=5$ is one percent of the mass, not a fifth.

More than those pairs is true. Sweeping the planning aggregate over $\hat K/K^*\in[0.5,2]$ at
fourteen points, $\Phi(K,p(\hat K))$ is strictly \emph{decreasing} in $\hat K$ at every one of the
thirteen steps, for $\gamma=1,3,5$ and for each of $K=0.90K^*,K^*,1.10K^*$. So \eqref{eq:W} holds
not merely at the pairs one happens to test but as a monotone comparative static in the plan. Over
all $91$ pairs at five pricing points each, $\gamma=3$ and $\gamma=5$ give no failures at all; the
$42$ failures out of $455$ at $\gamma=1$ all involve the extreme candidate $\hat K=0.5K^*$, whose
carried aggregate is $48.6$, ten times the capital stock and far outside the bracket of
Theorem~\ref{thm:exist5}.

